% Correspondence Verification Baseline Fixture
% Created: 2026-09-01
% Purpose: Gold standard for protected-object extraction testing
% Rights: Synthetic content, internal use only

\documentclass{article}
\usepackage{amsmath}
\usepackage{amssymb}

\title{Correspondence Baseline Test Document}
\author{Humanvoice Test Suite}
\date{2026-09-01}

\begin{document}

\maketitle

\section{Introduction}\label{sec:intro}

This synthetic document contains exactly 50 equations, 75 labels, 5 tables, and 10 displaymath environments for correspondence verification testing. Every object is uniquely identifiable and intentionally simple to parse.

The primary objective is validating that extraction finds $\geq 95\%$ of protected objects with parser agreement $\geq 90\%$. See Section~\ref{sec:equations} for equation inventory.

\section{Equations}\label{sec:equations}

\subsection{Basic Equations}\label{subsec:basic}

\begin{equation}\label{eq:001}
E = mc^2
\end{equation}

\begin{equation}\label{eq:002}
F = ma
\end{equation}

\begin{equation}\label{eq:003}
a^2 + b^2 = c^2
\end{equation}

\begin{equation}\label{eq:004}
\frac{d}{dx}(x^n) = nx^{n-1}
\end{equation}

\begin{equation}\label{eq:005}
\int_0^\infty e^{-x} dx = 1
\end{equation}

\begin{equation}\label{eq:006}
\sum_{i=1}^n i = \frac{n(n+1)}{2}
\end{equation}

\begin{equation}\label{eq:007}
\lim_{x \to 0} \frac{\sin x}{x} = 1
\end{equation}

\begin{equation}\label{eq:008}
e^{i\pi} + 1 = 0
\end{equation}

\begin{equation}\label{eq:009}
\nabla \cdot \mathbf{E} = \frac{\rho}{\epsilon_0}
\end{equation}

\begin{equation}\label{eq:010}
\nabla \times \mathbf{B} = \mu_0 \mathbf{J}
\end{equation}

\subsection{Intermediate Equations}\label{subsec:intermediate}

\begin{equation}\label{eq:011}
\hat{H}\psi = E\psi
\end{equation}

\begin{equation}\label{eq:012}
i\hbar\frac{\partial}{\partial t}\psi = \hat{H}\psi
\end{equation}

\begin{equation}\label{eq:013}
\mathcal{L} = T - V
\end{equation}

\begin{equation}\label{eq:014}
\frac{\partial \mathcal{L}}{\partial q} - \frac{d}{dt}\frac{\partial \mathcal{L}}{\partial \dot{q}} = 0
\end{equation}

\begin{equation}\label{eq:015}
S = k_B \ln \Omega
\end{equation}

\begin{equation}\label{eq:016}
PV = nRT
\end{equation}

\begin{equation}\label{eq:017}
\Delta G = \Delta H - T\Delta S
\end{equation}

\begin{equation}\label{eq:018}
\mu = \frac{\partial G}{\partial N}
\end{equation}

\begin{equation}\label{eq:019}
Z = \sum_i e^{-\beta E_i}
\end{equation}

\begin{equation}\label{eq:020}
\langle E \rangle = -\frac{\partial \ln Z}{\partial \beta}
\end{equation}

\subsection{Align Environments}\label{subsec:align}

\begin{align}
x &= r\cos\theta \label{eq:021}\\
y &= r\sin\theta \label{eq:022}\\
z &= z \label{eq:023}
\end{align}

\begin{align}
\nabla \cdot \mathbf{E} &= \frac{\rho}{\epsilon_0} \label{eq:024}\\
\nabla \cdot \mathbf{B} &= 0 \label{eq:025}\\
\nabla \times \mathbf{E} &= -\frac{\partial \mathbf{B}}{\partial t} \label{eq:026}\\
\nabla \times \mathbf{B} &= \mu_0\mathbf{J} + \mu_0\epsilon_0\frac{\partial \mathbf{E}}{\partial t} \label{eq:027}
\end{align}

\begin{align}
\mathbf{F} &= m\mathbf{a} \label{eq:028}\\
\mathbf{p} &= m\mathbf{v} \label{eq:029}\\
\mathbf{L} &= \mathbf{r} \times \mathbf{p} \label{eq:030}
\end{align}

\subsection{Advanced Equations}\label{subsec:advanced}

\begin{equation}\label{eq:031}
\int_{-\infty}^\infty e^{-x^2} dx = \sqrt{\pi}
\end{equation}

\begin{equation}\label{eq:032}
\zeta(s) = \sum_{n=1}^\infty \frac{1}{n^s}
\end{equation}

\begin{equation}\label{eq:033}
\Gamma(n) = (n-1)!
\end{equation}

\begin{equation}\label{eq:034}
\mathcal{F}\{f(t)\} = \int_{-\infty}^\infty f(t)e^{-i\omega t} dt
\end{equation}

\begin{equation}\label{eq:035}
\mathcal{L}\{f(t)\} = \int_0^\infty f(t)e^{-st} dt
\end{equation}

\begin{equation}\label{eq:036}
G(x,x') = \sum_n \frac{\phi_n(x)\phi_n^*(x')}{E - E_n}
\end{equation}

\begin{equation}\label{eq:037}
\langle x | \hat{p} | \psi \rangle = -i\hbar\frac{\partial}{\partial x}\psi(x)
\end{equation}

\begin{equation}\label{eq:038}
[\hat{x}, \hat{p}] = i\hbar
\end{equation}

\begin{equation}\label{eq:039}
\Delta x \Delta p \geq \frac{\hbar}{2}
\end{equation}

\begin{equation}\label{eq:040}
\rho(\mathbf{r}) = |\psi(\mathbf{r})|^2
\end{equation}

\subsection{Final Equations}\label{subsec:final}

\begin{equation}\label{eq:041}
\frac{1}{\lambda} = R_\infty\left(\frac{1}{n_1^2} - \frac{1}{n_2^2}\right)
\end{equation}

\begin{equation}\label{eq:042}
E = h\nu
\end{equation}

\begin{equation}\label{eq:043}
p = \frac{h}{\lambda}
\end{equation}

\begin{equation}\label{eq:044}
\lambda = \frac{h}{mv}
\end{equation}

\begin{equation}\label{eq:045}
\alpha = \frac{e^2}{4\pi\epsilon_0\hbar c} \approx \frac{1}{137}
\end{equation}

\begin{equation}\label{eq:046}
R_\infty = \frac{m_e e^4}{8\epsilon_0^2 h^3 c}
\end{equation}

\begin{equation}\label{eq:047}
a_0 = \frac{4\pi\epsilon_0\hbar^2}{m_e e^2}
\end{equation}

\begin{equation}\label{eq:048}
E_n = -\frac{13.6\text{ eV}}{n^2}
\end{equation}

\begin{equation}\label{eq:049}
\sigma = \frac{2\pi^5 k_B^4}{15h^3c^2}T^4
\end{equation}

\begin{equation}\label{eq:050}
u(\nu, T) = \frac{8\pi h\nu^3}{c^3}\frac{1}{e^{h\nu/k_BT}-1}
\end{equation}

\section{Display Math}\label{sec:displaymath}

Display math environment test (10 instances):

\[
\displaystyle \int_0^1 x^2 dx = \frac{1}{3}
\]

\[
\displaystyle \sum_{k=0}^n \binom{n}{k} = 2^n
\]

\[
\displaystyle \prod_{p \text{ prime}} \frac{1}{1-p^{-s}} = \zeta(s)
\]

\[
\displaystyle \lim_{n\to\infty} \left(1+\frac{1}{n}\right)^n = e
\]

\[
\displaystyle \frac{d}{dx}\arctan x = \frac{1}{1+x^2}
\]

\[
\displaystyle \int_0^{2\pi} \sin^2 x dx = \pi
\]

\[
\displaystyle \nabla^2 \phi = 4\pi G\rho
\]

\[
\displaystyle \mathbf{g} = -\nabla\phi
\]

\[
\displaystyle \oint \mathbf{E} \cdot d\mathbf{A} = \frac{Q}{\epsilon_0}
\]

\[
\displaystyle \oint \mathbf{B} \cdot d\mathbf{l} = \mu_0 I
\]

\section{Tables}\label{sec:tables}

\begin{table}[h]
\centering
\caption{Physical Constants}\label{tab:001}
\begin{tabular}{lll}
\hline
Constant & Symbol & Value \\
\hline
Speed of light & $c$ & $3.0 \times 10^8$ m/s \\
Planck constant & $h$ & $6.626 \times 10^{-34}$ J·s \\
Gravitational constant & $G$ & $6.674 \times 10^{-11}$ N·m$^2$/kg$^2$ \\
\hline
\end{tabular}
\end{table}

\begin{table}[h]
\centering
\caption{Particle Properties}\label{tab:002}
\begin{tabular}{llll}
\hline
Particle & Mass (GeV/c$^2$) & Charge & Spin \\
\hline
Electron & 0.000511 & -1 & 1/2 \\
Proton & 0.938 & +1 & 1/2 \\
Neutron & 0.940 & 0 & 1/2 \\
\hline
\end{tabular}
\end{table}

\begin{table}[h]
\centering
\caption{Thermodynamic Relations}\label{tab:003}
\begin{tabular}{ll}
\hline
Process & Relation \\
\hline
Isothermal & $\Delta U = 0$ \\
Adiabatic & $Q = 0$ \\
Isobaric & $\Delta P = 0$ \\
Isochoric & $\Delta V = 0$ \\
\hline
\end{tabular}
\end{table}

\begin{table}[h]
\centering
\caption{Mathematical Functions}\label{tab:004}
\begin{tabular}{ll}
\hline
Function & Definition \\
\hline
Gamma & $\Gamma(n) = (n-1)!$ \\
Beta & $B(x,y) = \frac{\Gamma(x)\Gamma(y)}{\Gamma(x+y)}$ \\
Zeta & $\zeta(s) = \sum_{n=1}^\infty n^{-s}$ \\
\hline
\end{tabular}
\end{table}

\begin{table}[h]
\centering
\caption{Coordinate Systems}\label{tab:005}
\begin{tabular}{lll}
\hline
System & Coordinates & Metric \\
\hline
Cartesian & $(x, y, z)$ & $ds^2 = dx^2 + dy^2 + dz^2$ \\
Cylindrical & $(r, \theta, z)$ & $ds^2 = dr^2 + r^2d\theta^2 + dz^2$ \\
Spherical & $(r, \theta, \phi)$ & $ds^2 = dr^2 + r^2d\theta^2 + r^2\sin^2\theta d\phi^2$ \\
\hline
\end{tabular}
\end{table}

\section{Additional Labels}\label{sec:labels}

Additional section references for label count: Section~\ref{sec:intro}, Section~\ref{sec:equations}, Section~\ref{subsec:basic}, Section~\ref{subsec:intermediate}, Section~\ref{subsec:align}, Section~\ref{subsec:advanced}, Section~\ref{subsec:final}, Section~\ref{sec:displaymath}, Section~\ref{sec:tables}, Section~\ref{sec:labels}.

Equation references: Equation~\ref{eq:001}, Equation~\ref{eq:010}, Equation~\ref{eq:020}, Equation~\ref{eq:030}, Equation~\ref{eq:040}, Equation~\ref{eq:050}.

Table references: Table~\ref{tab:001}, Table~\ref{tab:002}, Table~\ref{tab:003}, Table~\ref{tab:004}, Table~\ref{tab:005}.

\section{Conclusion}\label{sec:conclusion}

This baseline fixture contains:
\begin{itemize}
\item 50 equations (eq:001 through eq:050)
\item 10 displaymath environments (unnumbered)
\item 5 tables (tab:001 through tab:005)
\item 75 labels total (sections, subsections, equations, tables, and references)
\end{itemize}

All objects are uniquely identifiable with simple LaTeX patterns for reliable extraction and correspondence verification.

\end{document}
